\section{Dữ liệu và chống rò rỉ}
\label{sec:data}

\subsection{Visual Genome}
Visual Genome v1.2~\cite{krishna2017vg} (giấy phép CC BY 4.0) có 108.077 ảnh và 5.408.689 câu mô tả vùng (region description). Với mỗi ảnh, nhóm lấy mẫu 5 câu mô tả vùng (riêng một ảnh chỉ có 3 câu), tổng cộng 540.383 câu. Ảnh và câu được mã hóa một lần bằng OpenCLIP ViT-B/32 (\code{laion2b\_s34b\_b79k}) trên Kaggle, lưu dưới dạng \code{.npz} để cả nhóm dùng chung. Phần dựng đồ thị dùng thêm tệp \code{objects.json} v1.4 với 2.516.939 đối tượng, trong đó 19,6\% không gắn được synset WordNet~\cite{miller1995wordnet} và 1.793 ảnh không có đối tượng nào. Thông tin nguồn, giấy phép và mã sha256 của từng tệp được ghi trong data card \code{data\_cards/visual\_genome.md}.

\subsection{Chia tập cố định và xử lý ảnh trùng}
Visual Genome không có cách chia chính thức. Nhóm chia theo \emph{ảnh}: 5.000 ảnh test, 5.000 ảnh val và 98.077 ảnh train; mỗi câu thuộc về tập của ảnh tương ứng. Dữ liệu có ảnh trùng: 5 cặp ảnh có embedding giống hệt nhau và 247 cặp có cosine $\ge 0,99$ (ví dụ ảnh \code{285982} và \code{2375578} có cosine 0,99992). Nếu hai bản của cùng một ảnh rơi vào train và test, dữ liệu sẽ bị rò rỉ. Vì vậy các ảnh có cosine $\ge 0,99$ được gom thành nhóm theo quan hệ bắc cầu (A gần trùng B, B gần trùng C thì cả ba chung nhóm) và cả nhóm luôn nằm trong cùng một tập. Ngưỡng này được chốt trước khi có bất kỳ kết quả truy vấn nào. Tệp chia tập \code{split\_seed0.json} được commit cùng mã sha256 và là nguồn duy nhất cho mọi thí nghiệm; chạy lại script cho ra tệp giống hệt từng byte.

\subsection{Giao thức đánh giá}
\label{sec:protocol}
Truy vấn là toàn bộ câu mô tả vùng đã lấy mẫu của ảnh test (25.000 câu), kho tìm kiếm là 5.000 ảnh test. Một ảnh được coi là đáp án đúng nếu nó có một câu trùng với truy vấn (cosine văn bản $\ge 0,9999$, ngưỡng này chỉ để bù sai số số học), kể cả ảnh gốc của câu. Lý do là 26\% số câu trong dữ liệu bị trùng, ví dụ ``a man'' hay ``white clouds''; nếu chỉ tính ảnh gốc là đúng thì CLIP bị đánh giá thấp một cách sai lệch. Trên tập test, 11,28\% truy vấn có nhiều hơn một đáp án.

Các chỉ số dùng trong báo cáo là Recall@1/5/10, \mrr{} (ảnh đúng nằm ngoài 100 kết quả đầu được tính bằng 0) và nDCG@10 nhị phân~\cite{jarvelin2002ndcg}. Vì mỗi ảnh chỉ xét 5 câu, Recall tuyệt đối là một cận dưới; so sánh giữa các mô hình vẫn công bằng vì mọi mô hình dùng chung định nghĩa. CLIP thuần và CLIP kết hợp đồ thị gọi cùng một hàm lấy 100 ứng viên, nên hai mô hình chỉ khác nhau ở thứ tự sắp xếp bên trong 100 ảnh đó.

\subsection{Các quy tắc chống rò rỉ}
Nhóm áp dụng năm quy tắc, đều được ghi thành ADR trước khi chạy thí nghiệm:
\begin{enumerate}
\item Từ vựng khái niệm, từ điển tên đối tượng $\to$ synset và cạnh đồng xuất hiện chỉ được đếm trên tập train. Một bài kiểm thử tự động xác nhận rằng ảnh test không làm thay đổi từ vựng.
\item Cạnh của ảnh val và test chỉ phụ thuộc vào vectơ CLIP của chính ảnh đó, không dùng nhãn (Mục~\ref{sec:graph}).
\item Mọi lựa chọn (learning rate, epoch, $\alpha$, phương pháp giải thích, số cạnh $m$) đều làm trên val. Tập test được chạy đúng một lần cho mỗi seed; script từ chối lần chạy thứ hai của cùng một thí nghiệm và seed.
\item Quy tắc chọn cấu hình chính và quy tắc của các ablation được commit trước khi chạy test. Kết quả âm vẫn được báo cáo.
\item Seed 0, 1, 2 chỉ thay đổi phần huấn luyện (thứ tự batch, mẫu âm khó); cách chia tập giữ nguyên. Vì CLIP được đóng băng và FAISS Flat là tất định, kết quả của CLIP thuần có độ lệch chuẩn bằng 0.
\end{enumerate}

\subsection{Flickr30k}
Tập công khai thứ hai là Flickr30k~\cite{young2014flickr} theo bản phân chia của Karpathy~\cite{karpathy2015}, gồm 31.014 ảnh, mỗi ảnh có 5 chú thích. Embedding được trích bằng cùng mô hình CLIP như Visual Genome. Nhóm chia lại tập theo đúng quy trình trên (21.014 ảnh train, 5.000 val, 5.000 test, seed 0; có 2 cặp ảnh cosine $\ge 0,99$ được gom chung) và không dùng trường \code{split} có sẵn của bản Karpathy. Truy vấn là cả 5 chú thích của mỗi ảnh val/test (25.000 câu); chỉ 0,08\% truy vấn val có hơn một đáp án.

Flickr30k không có annotation đối tượng nên không đếm được từ vựng khái niệm như trên Visual Genome. Nhóm chép nguyên 3.880 đỉnh khái niệm (synset và vectơ văn bản CLIP) từ đồ thị Visual Genome, rồi nối mỗi ảnh Flickr30k với 16 khái niệm gần nhất theo cosine CLIP, được tổng cộng 496.224 cạnh. Vì từ vựng chỉ được đếm trên train Visual Genome và cạnh của mỗi ảnh chỉ phụ thuộc vào vectơ của chính ảnh đó, quy trình này không làm rò rỉ thông tin từ test Flickr30k. Mô hình có cạnh đồng xuất hiện (cần annotation đối tượng) và phần phân nhóm truy vấn (luật được chốt trên câu Visual Genome) không được chạy trên Flickr30k.

\subsection{MMKG tiếng Việt tự xây}
\label{sec:mmkg-data}
Nhóm chọn miền hẹp là \emph{di tích và di sản Việt Nam}: các thực thể có danh hiệu di sản trên Wikidata (\code{P1435}) và thuộc Việt Nam (\code{P17 = Q881}). Thực thể, bộ ba và ảnh được thu tự động từ Wikidata và Wikimedia Commons. Dữ liệu thô (snapshot) được lưu nguyên trạng rồi mới lọc ở bước dựng, nên mọi bước lọc đều kiểm tra lại được. Quy trình gồm ba bước:
\begin{enumerate}
\item \textbf{Thu thập}: lấy mọi phát biểu của thực thể trong miền trỏ tới một thực thể Wikidata khác. Dữ liệu thô có 848 thực thể, 3.581 bộ ba và 5.093 cặp (thực thể, ảnh) ứng viên.
\item \textbf{Dựng}: chỉ giữ 18 thuộc tính Wikidata, đổi thành 17 quan hệ như \code{heritage\_designation}, \code{instance\_of}, \code{kept\_at}, \code{made\_from}, \code{located\_in}, \code{culture}. Ảnh phải là \code{P18} của thực thể hoặc nằm trực tiếp trong danh mục Commons của nó; ảnh thuộc nhiều thực thể, ảnh không phải dạng raster, ảnh có cạnh ngắn dưới 224 px và ảnh có giấy phép không mở đều bị loại. Mỗi hiện vật giữ từ 3 đến 5 ảnh.
\item \textbf{Duyệt}: toàn bộ 962 dòng được một mô hình AI duyệt sơ bộ, sau đó một thành viên duyệt lại từng dòng. Ảnh có mặt người, ảnh không đúng hiện vật và ảnh trùng bị loại.
\end{enumerate}
Bộ chính thức có 119 thực thể (80 hiện vật có ảnh và 39 khái niệm, tất cả đều có tên tiếng Việt), 303 bộ ba và 380 ảnh. Mọi ảnh có giấy phép tự do (CC0, public domain, CC BY, CC BY-SA), kèm tác giả và đường dẫn nguồn. Các quan hệ nhiều nhất là \code{heritage\_designation} (86), \code{instance\_of} (65), \code{kept\_at} (45) và \code{made\_from} (30). Thông tin đầy đủ về nguồn, giấy phép, cách tạo bộ ba và hạn chế của dữ liệu nằm trong \code{data\_cards/mmkg\_vi.md}.

\begin{figure}[t]
\centering
\begin{tikzpicture}[
  ent/.style={draw=blue!60!black,thick,rounded corners=2pt,fill=blue!6,font=\scriptsize,align=center,text width=2.0cm,minimum height=0.6cm},
  con/.style={draw=orange!70!black,rounded corners=6pt,fill=orange!8,font=\scriptsize,align=center,text width=1.9cm,minimum height=0.5cm},
  pic/.style={draw=gray,fill=gray!10,font=\tiny,minimum width=0.6cm,minimum height=0.45cm},
  r/.style={-{Stealth[length=1.4mm]},thick,blue!50!black},
  lab/.style={font=\tiny,fill=white,inner sep=1pt}]
\node[ent] (h) at (0,0) {Hiện vật\\(\code{item})};
\node[con] (t1) at (-2.4,1.5) {Bảo vật quốc gia};
\node[con] (t2) at (2.4,1.5) {Văn hóa Đông Sơn};
\node[con] (t3) at (-2.4,-1.4) {Đồng thanh};
\node[con] (t4) at (2.4,-1.4) {Bảo tàng};
\draw[r] (h) -- node[lab]{\code{heritage\_designation}} (t1);
\draw[r] (h) -- node[lab]{\code{culture}} (t2);
\draw[r] (h) -- node[lab]{\code{made\_from}} (t3);
\draw[r] (h) -- node[lab]{\code{kept\_at}} (t4);
\node[pic] (p1) at (-0.75,-2.3) {ảnh};
\node[pic] (p2) at (0,-2.3) {ảnh};
\node[pic] (p3) at (0.75,-2.3) {ảnh};
\draw[gray,thick] (h) -- (p1); \draw[gray,thick] (h) -- (p2); \draw[gray,thick] (h) -- (p3);
\end{tikzpicture}
\caption{Lược đồ MMKG tiếng Việt: mỗi hiện vật nối với các khái niệm qua 17 loại quan hệ lấy từ Wikidata và có 3--5 ảnh từ Wikimedia Commons. Câu truy vấn gián tiếp được sinh từ các bộ ba này.}
\label{fig:mmkg}
\end{figure}

Hình~\ref{fig:mmkg} mô tả lược đồ của MMKG. Từ MMKG, nhóm sinh tự động 318 truy vấn. Mỗi hiện vật có một câu \emph{trực tiếp} gọi tên hiện vật và một câu \emph{gián tiếp} giấu tên, chỉ mô tả hiện vật qua các bộ ba của nó (ví dụ ``a heritage of Vietnam, made of bronze, of the Dong Son culture, kept at National Museum of Vietnamese History''). Mỗi câu có một bản tiếng Anh và một bản tiếng Việt cùng nội dung. Đáp án là mọi ảnh của mọi hiện vật khớp với câu. Câu gián tiếp được chia thành loại \emph{cụ thể} (chỉ một hiện vật khớp, 35 câu mỗi ngôn ngữ) và loại \emph{chung chung} (nhiều hiện vật khớp, 44 câu). Một hiện vật không có bộ ba nào nên không có câu gián tiếp. MMKG chỉ được dùng để đánh giá; không có lựa chọn nào được thực hiện trên bộ này.
